% GENERATED FILE. Edit canonical lessons, then run npm run generate:books.
% canonical-source-hash: fnv1a64:0af786b3d5c27fc4
% canonical-lessons: ML-C208-accar, ML-C208-moru, ML-C208-neyyu, ML-C208-venna, ML-C208-ten

\chapter{Pickle, Buttermilk, Ghee, Butter, Honey}
\label{ch:ml-pickle-buttermilk-ghee-butter-honey}
\hlchaptermodality{208}

\begin{chapteropening}
\textbf{By the end of this chapter:} \emph{I can say a pickle, buttermilk, ghee, butter and honey.}

Complete the last lesson of chapter 208: I can say a pickle, buttermilk, ghee, butter and honey.
\end{chapteropening}

\section[\texorpdfstring{\ml{അച്ചാർ} (accār)}{അച്ചാർ (accār)}]{\ml{അച്ചാർ} (accār) — a pickle}

\label{lesson:ML-C208-accar}



\begin{quote}
Before the new one: say the Malayalam for a ladle, then the Malayalam for a sieve.
\end{quote}

\subsection*{You'll want to know: \ml{അച്ചാർ}}
\textbf{\ml{അച്ചാർ}} — \emph{accār} — \textquotedblleft{}a pickle\textquotedblright{}.

\begin{cousinweb}
\textbf{\ml{അച്ചാർ}} (\emph{accār}) $\leftarrow$ Persian \emph{āchār}, \textquotedblleft{}pickle\textquotedblright{}, which came in through Hindustani.
\end{cousinweb}

\begin{grammarlens}[title={What you've built}]
One more word for this chapter.
\end{grammarlens}

\subsection*{Your turn}
\begin{itemize}
\raggedright
  \item \emph{Say it:} \emph{accār}
  \item \emph{Say it:} \emph{accār}, once more
  \item \emph{Say it:} say \emph{accār}
  \item \emph{Recall:} say the Malayalam for a ladle, then the Malayalam for a sieve, then say \emph{accār} again
\end{itemize}

\begin{tcolorbox}[breakable,colback=teal!4,colframe=teal!35!black,title={Before you move on}]
What does \ml{അച്ചാർ} mean? (\textquotedblleft{}A pickle\textquotedblright{}.) Say it once more.
\end{tcolorbox}
\section[\texorpdfstring{\ml{മോര്} (mōrŭ)}{മോര് (mōrŭ)}]{\ml{മോര്} (mōrŭ) — buttermilk}

\label{lesson:ML-C208-moru}



\begin{quote}
Before the new one: say the Malayalam for a sieve, then the Malayalam for a pickle.
\end{quote}

\subsection*{You'll want to know: \ml{മോര്}}
\textbf{\ml{മോര്}} — \emph{mōrŭ} — \textquotedblleft{}buttermilk\textquotedblright{}.

\begin{cousinweb}
\textbf{\ml{മോര്}} (\emph{mōrŭ}) is native Dravidian, one word with Tamil \emph{mucar} and Kannada \emph{mosar(u)} (DEDR 4902). DEDR 4902 compares Sanskrit \emph{moraṭa}.
\end{cousinweb}

\begin{grammarlens}[title={What you've built}]
One more word for this chapter.
\end{grammarlens}

\subsection*{Your turn}
\begin{itemize}
\raggedright
  \item \emph{Say it:} \emph{mōrŭ}
  \item \emph{Say it:} \emph{mōrŭ}, once more
  \item \emph{Say it:} say \emph{mōrŭ}
  \item \emph{Recall:} say the Malayalam for a sieve, then the Malayalam for a pickle, then say \emph{mōrŭ} again
\end{itemize}

\begin{tcolorbox}[breakable,colback=teal!4,colframe=teal!35!black,title={Before you move on}]
What does \ml{മോര്} mean? (\textquotedblleft{}Buttermilk\textquotedblright{}.) Say it once more.
\end{tcolorbox}
\section[\texorpdfstring{\ml{നെയ്യ്} (neyyŭ)}{നെയ്യ് (neyyŭ)}]{\ml{നെയ്യ്} (neyyŭ) — ghee}

\label{lesson:ML-C208-neyyu}



\begin{quote}
Before the new one: say the Malayalam for a pickle, then the Malayalam for buttermilk.
\end{quote}

\subsection*{You'll want to know: \ml{നെയ്യ്}}
\textbf{\ml{നെയ്യ്}} — \emph{neyyŭ} — \textquotedblleft{}ghee\textquotedblright{}.

\begin{cousinweb}
\textbf{\ml{നെയ്യ്}} (\emph{neyyŭ}) is native Dravidian, one word with Tamil \emph{ney}, Kannada \emph{ney} and Telugu \emph{neyyi} (DEDR 3746). Its older sense is any fat or oil.
\end{cousinweb}

\begin{grammarlens}[title={What you've built}]
One more word for this chapter.
\end{grammarlens}

\subsection*{Your turn}
\begin{itemize}
\raggedright
  \item \emph{Say it:} \emph{neyyŭ}
  \item \emph{Say it:} \emph{neyyŭ}, once more
  \item \emph{Say it:} say \emph{neyyŭ}
  \item \emph{Recall:} say the Malayalam for a pickle, then the Malayalam for buttermilk, then say \emph{neyyŭ} again
\end{itemize}

\begin{tcolorbox}[breakable,colback=teal!4,colframe=teal!35!black,title={Before you move on}]
What does \ml{നെയ്യ്} mean? (\textquotedblleft{}Ghee\textquotedblright{}.) Say it once more.
\end{tcolorbox}
\section[\texorpdfstring{\ml{വെണ്ണ} (veṇṇa)}{വെണ്ണ (veṇṇa)}]{\ml{വെണ്ണ} (veṇṇa) — butter}

\label{lesson:ML-C208-venna}



\begin{quote}
Before the new one: say the Malayalam for buttermilk, then the Malayalam for ghee.
\end{quote}

\subsection*{You'll want to know: \ml{വെണ്ണ}}
\textbf{\ml{വെണ്ണ}} — \emph{veṇṇa} — \textquotedblleft{}butter\textquotedblright{}.

\begin{cousinweb}
\textbf{\ml{വെണ്ണ}} (\emph{veṇṇa}) is native Dravidian, \textquotedblleft{}white fat\textquotedblright{}: DEDR 5496 lists Tamil \emph{veṇ-ṇey}, \emph{veṇṇai}, \textquotedblleft{}butter\textquotedblright{}, the word for \textquotedblleft{}white\textquotedblright{} joined to the word for fat.
\end{cousinweb}

\begin{grammarlens}[title={What you've built}]
One more word for this chapter.
\end{grammarlens}

\subsection*{Your turn}
\begin{itemize}
\raggedright
  \item \emph{Say it:} \emph{veṇṇa}
  \item \emph{Say it:} \emph{veṇṇa}, once more
  \item \emph{Say it:} say \emph{veṇṇa}
  \item \emph{Recall:} say the Malayalam for buttermilk, then the Malayalam for ghee, then say \emph{veṇṇa} again
\end{itemize}

\begin{tcolorbox}[breakable,colback=teal!4,colframe=teal!35!black,title={Before you move on}]
What does \ml{വെണ്ണ} mean? (\textquotedblleft{}Butter\textquotedblright{}.) Say it once more.
\end{tcolorbox}
\section[\texorpdfstring{\ml{തേൻ} (tēn)}{തേൻ (tēn)}]{\ml{തേൻ} (tēn) — honey}

\label{lesson:ML-C208-ten}



\begin{quote}
Before the new one: say the Malayalam for ghee, then the Malayalam for butter.
\end{quote}

\subsection*{You'll want to know: \ml{തേൻ}}
\textbf{\ml{തേൻ}} — \emph{tēn} — \textquotedblleft{}honey\textquotedblright{}.

\begin{cousinweb}
\textbf{\ml{തേൻ}} (\emph{tēn}) is native Dravidian, one word with Tamil \emph{tēṉ} and Kannada \emph{jēnu}, \textquotedblleft{}honey\textquotedblright{} (DEDR 3268). DEDR 3268 compares, with a query, Prakrit \emph{tiṇisa}, \textquotedblleft{}honeycomb\textquotedblright{}.
\end{cousinweb}

\begin{grammarlens}[title={What you've built}]
One more word for this chapter.
\end{grammarlens}

\subsection*{Your turn}
\begin{itemize}
\raggedright
  \item \emph{Say it:} \emph{tēn}
  \item \emph{Say it:} \emph{tēn}, once more
  \item \emph{Say it:} say \emph{tēn}
  \item \emph{Recall:} say the Malayalam for ghee, then the Malayalam for butter, then say \emph{tēn} again
\end{itemize}

\begin{tcolorbox}[breakable,colback=teal!4,colframe=teal!35!black,title={Before you move on}]
What does \ml{തേൻ} mean? (\textquotedblleft{}Honey\textquotedblright{}.) Say it once more.
\end{tcolorbox}
